\textbf{Phần 1. Trắc nghiệm lựa chọn}
\setcounterpageref{ex}{0}

\begin{ex} %[3D2N11]
Kết quả của phép nhân $8 \times 4$ là
\choice
{\True $32$}
{$28$}
{$36$}
{$40$}
\loigiai{
$8 \times 4 = 32$.\\
\textbf{Đáp án của câu hỏi là A}
}
\end{ex}

\begin{ex} %[3D2N11]
Kết quả của phép chia $40 : 8$ là
\choice
{$4$}
{$6$}
{\True $5$}
{$7$}
\loigiai{
$40 : 8 = 5$.\\
\textbf{Đáp án của câu hỏi là C}
}
\end{ex}

\begin{ex} %[3D2H11]
Kết quả của phép nhân $8 \times 7$ là
\choice
{\True $56$}
{$54$}
{$58$}
{$64$}
\loigiai{
$8 \times 7 = 56$.\\
\textbf{Đáp án của câu hỏi là A}
}
\end{ex}

\begin{ex} %[3D2H11]
Kết quả của phép chia $64 : 8$ là
\choice
{$7$}
{\True $8$}
{$9$}
{$6$}
\loigiai{
$64 : 8 = 8$.\\
\textbf{Đáp án của câu hỏi là B}
}
\end{ex}

\begin{ex} %[3D2H11]
Điền số thích hợp vào chỗ chấm: $8 \times \dots = 72$.
\choice
{$8$}
{$7$}
{\True $9$}
{$6$}
\loigiai{
$72 : 8 = 9$.\\
\textbf{Đáp án của câu hỏi là C}
}
\end{ex}

\begin{ex} %[3D2H11]
Điền số thích hợp vào chỗ chấm: $\dots : 8 = 6$.
\choice
{$40$}
{$56$}
{$36$}
{\True $48$}
\loigiai{
$6 \times 8 = 48$.\\
\textbf{Đáp án của câu hỏi là D}
}
\end{ex}

\begin{ex} %[3D2V11]
Tính giá trị biểu thức: $8 \times 5 + 35$.
\choice
{$70$}
{\True $75$}
{$80$}
{$65$}
\loigiai{
$8 \times 5 + 35 = 40 + 35 = 75$.\\
\textbf{Đáp án của câu hỏi là B}
}
\end{ex}

\begin{ex} %[3D2V11]
Một con nhện có 8 cái chân. Hỏi 6 con nhện có tất cả bao nhiêu cái chân?
\choice
{$42$ cái chân}
{$54$ cái chân}
{$40$ cái chân}
{\True $48$ cái chân}
\loigiai{
$8 \times 6 = 48$ cái chân.\\
\textbf{Đáp án của câu hỏi là D}
}
\end{ex}

\begin{ex} %[3D2V11]
Có 72 cái bánh chia đều vào 8 hộp. Mỗi hộp có số cái bánh là
\choice
{\True $9$ cái}
{$8$ cái}
{$7$ cái}
{$10$ cái}
\loigiai{
$72 : 8 = 9$ cái bánh.\\
\textbf{Đáp án của câu hỏi là A}
}
\end{ex}

\begin{ex} %[3D2V11]
Một bao gạo có $56\text{ kg}$ chia đều vào 8 túi. Mỗi túi có số ki-lô-gam gạo là
\choice
{$6\text{ kg}$}
{$8\text{ kg}$}
{\True $7\text{ kg}$}
{$9\text{ kg}$}
\loigiai{
$56 : 8 = 7\text{ kg}$.\\
\textbf{Đáp án của câu hỏi là C}
}
\end{ex}

\begin{ex} %[3D2V11]
Một hình chữ nhật có chiều rộng $8\text{ cm}$, chiều dài gấp 3 lần chiều rộng. Chiều dài hình chữ nhật là
\choice
{$20\text{ cm}$}
{$28\text{ cm}$}
{$16\text{ cm}$}
{\True $24\text{ cm}$}
\loigiai{
$8 \times 3 = 24\text{ cm}$.\\
\textbf{Đáp án của câu hỏi là D}
}
\end{ex}

\begin{ex} %[3D2C11]
Có 8 thùng dầu, mỗi thùng chứa $9\text{ l}$ dầu. Người ta đã lấy ra $25\text{ l}$ dầu. Số dầu còn lại là
\choice
{$45\text{ l}$}
{\True $47\text{ l}$}
{$49\text{ l}$}
{$52\text{ l}$}
\loigiai{
Tổng số dầu: $8 \times 9 = 72\text{ l}$.\\
Còn lại: $72 - 25 = 47\text{ l}$.\\
\textbf{Đáp án của câu hỏi là B}
}
\end{ex}

\textbf{Phần 2. Trắc nghiệm đúng sai}
\setcounterpageref{ex}{0}

\begin{ex} %[3D2H11]
Xét tính đúng sai của các phép tính trong bảng nhân 8, bảng chia 8:
\choiceTF
{\True $8 \times 3 = 24$.}
{\True $32 : 8 = 4$.}
{$8 \times 8 = 62$.}
{\True $8 \times 8 = 64$.}
\loigiai{
\begin{itemize}
    \item $8 \times 3 = 24$: Đúng.
    \item $32 : 8 = 4$: Đúng.
    \item $8 \times 8 = 64 \neq 62$. Ý 3 Sai, Ý 4 Đúng.
\end{itemize}
\begin{itemize}
    \item a) \textbf{Đúng}.
    \item b) \textbf{Đúng}.
    \item c) \textbf{Sai}.
    \item d) \textbf{Đúng}.
\end{itemize}
}
\end{ex}

\begin{ex} %[3D2V11]
Có 48 học sinh chia đều thành 8 nhóm. Xét tính đúng sai:
\choiceTF
{\True Mỗi nhóm có 6 học sinh.}
{\True 3 nhóm có tất cả 18 học sinh.}
{5 nhóm có tất cả 35 học sinh.}
{\True 5 nhóm có tất cả 30 học sinh.}
\loigiai{
\begin{itemize}
    \item Mỗi nhóm: $48 : 8 = 6$ bạn: Đúng.
    \item 3 nhóm: $6 \times 3 = 18$ bạn: Đúng.
    \item 5 nhóm: $6 \times 5 = 30$ bạn $\neq 35$. Ý 3 Sai, Ý 4 Đúng.
\end{itemize}
\begin{itemize}
    \item a) \textbf{Đúng}.
    \item b) \textbf{Đúng}.
    \item c) \textbf{Sai}.
    \item d) \textbf{Đúng}.
\end{itemize}
}
\end{ex}

\textbf{Phần 3. Trả lời ngắn}
\setcounterpageref{ex}{0}

\begin{ex} %[3D2H11]
Tính tích: $8 \times 6$.
\shortans{48}
\loigiai{
\textbf{Đáp số:} $48$\\
$8 \times 6 = 48$.
}
\end{ex}

\begin{ex} %[3D2H11]
Tính thương: $24 : 8$.
\shortans{3}
\loigiai{
\textbf{Đáp số:} $3$\\
$24 : 8 = 3$.
}
\end{ex}

\begin{ex} %[3D2V11]
Tính giá trị biểu thức: $8 \times 9 - 35$.
\shortans{37}
\loigiai{
\textbf{Đáp số:} $37$\\
$72 - 35 = 37$.
}
\end{ex}

\begin{ex} %[3D2V11]
Một hộp có 8 gói bánh. Hỏi 9 hộp như thế có tất cả bao nhiêu gói bánh?
\shortans{72}
\loigiai{
\textbf{Đáp số:} $72$\\
$8 \times 9 = 72$ gói bánh.
}
\end{ex}

\begin{ex} %[3D2V11]
Có 80 mét vải may đều được 8 tấm rèm. Hỏi mỗi tấm rèm may hết bao nhiêu mét vải?
\shortans{10}
\loigiai{
\textbf{Đáp số:} $10$\\
$80 : 8 = 10\text{ m}$.
}
\end{ex}

\begin{ex} %[3D2C11]
Tìm $x$ biết: $x \times 8 = 50 + 14$.
\shortans{8}
\loigiai{
\textbf{Đáp số:} $8$\\
$x \times 8 = 64 \Rightarrow x = 64 : 8 = 8$.
}
\end{ex}
